\documentclass[10pt]{beamer}
\usepackage[utf8]{inputenc}
\usepackage[T1,T2A]{fontenc}
\usepackage[russian]{babel}
\usepackage{color}
\usepackage{calc}
\usepackage{graphicx}
\usepackage{epstopdf}
\usepackage{hyperref}
\hypersetup{unicode,colorlinks}
\usepackage{csquotes}
\usepackage{upquote}
\usepackage{cprotect}
\usetheme[progressbar=head,numbering=fraction,block=fill]{metropolis}
\usepackage{dejavu}
\usepackage{etoolbox}
\usepackage{bm}
\usepackage[ND]{prftree}
\usepackage[tableaux]{prooftrees}
\usepackage{mathtools} % for bigtimes
\usepackage{diagbox}
\usepackage{verbatimbox}
\usepackage{fitch}
% fitch по умолчанию печатает эти правила как $\Rightarrow$I/E; в курсе импликация всюду $\to$.
\gappto\ndrules{\def\ii{\by{$\to$I}}\def\ie{\by{$\to$E}}}
%\usepackage[defaultsans]{droidsans}
%\usepackage[defaultmono]{droidmono}
%\makeatletter
%\input droid/t1fdm.fd
%\makeatother

\makeatletter
\DeclareRobustCommand*{\bfseries}{%
    \not@math@alphabet\bfseries\mathbf
    \fontseries\bfdefault\selectfont
    \boldmath
}
\makeatother

\makeatletter
\chardef\straightquote@code=\catcode`'
\chardef\backquote@code=\catcode``
\catcode`'=\active \catcode``=\active
\patchcmd{\@noligs}
{\textasciigrave}
{\fixedtextasciigrave}
{}{}
\newcommand{\fixedtextasciigrave}{%
    \makebox[.5em]{\fontencoding{TS1}\fontfamily{fvs}\selectfont\textasciigrave}% Vera Sans
}
\catcode\lq\'=\straightquote@code
\catcode\lq\`=\backquote@code
\makeatother

\setbeamercovered{again covered={\opaqueness<1->{25}}}

\forestset{close with=\ensuremath{\bigtimes}}

\newcommand{\customframefont}[1]{
    \setbeamertemplate{itemize/enumerate body begin}{#1}
    \setbeamertemplate{itemize/enumerate subbody begin}{#1}
}

\NewEnviron{framefont}[1]{
    \customframefont{#1} % for itemize/enumerate
    {#1 % For the text outside itemize/enumerate
        \BODY
    }
    \customframefont{\normalsize}
}


\makeatletter
\DeclareRobustCommand{\rvdots}{%
    \vphantom{\int\limits^x}\smash[t]{\vdots}
}
\makeatother

\newcommand{\N}{\mathbb{N}}
\newcommand{\Z}{\mathbb{Z}}
\newcommand{\Q}{\mathbb{Q}}
\newcommand{\R}{\mathbb{R}}
\newcommand{\pow}[1]{\mathcal{P}({#1})}
\newcommand{\continuum}{\mathfrak{c}}

% Многобуквенные имена в формулах: курсив с нормальным кернингом (а не произведение букв).
\newcommand{\ident}[1]{\mathit{#1}}
\newcommand{\Var}{\ident{Var}}
\newcommand{\Prop}{\ident{Prop}}
\newcommand{\Const}{\ident{Const}}
\newcommand{\Fun}{\ident{Fun}}
\newcommand{\Pred}{\ident{Pred}}
\newcommand{\Term}{\ident{Term}}
\newcommand{\Frm}{\ident{Form}} % \Form занят hyperref
\newcommand{\arity}{\ident{arity}}
\newcommand{\Th}{\ident{Th}}
\newcommand{\Dom}{\ident{Dom}}
\newcommand{\Ran}{\ident{Ran}}
\newcommand{\id}{\ident{id}}
\newcommand{\Ord}{\ident{Ord}}
\newcommand{\RAA}{\ident{RAA}}
% Названия теорий и аксиом: прямой шрифт.
\newcommand{\ZF}{\mathrm{ZF}}
\newcommand{\ZFC}{\mathrm{ZFC}}
\newcommand{\AC}{\mathrm{AC}}
\newcommand{\DC}{\mathrm{DC}}
\newcommand{\CH}{\mathrm{CH}}
\newcommand{\GCH}{\mathrm{GCH}}

\author{Алексей Романов}
\subtitle{Математическая логика и теория алгоритмов}
%\logo{}
\institute{МИЭТ}
\subject{Математическая логика и теория алгоритмов}
%\setbeamercovered{transparent}
%\setbeamertemplate{navigation symbols}{}


\title{Деревья истинности\\для логики предикатов}
\date{\today}

\begin{document}
\begin{frame}[plain]
    \maketitle
\end{frame}

\begin{frame}
\frametitle{С чего начинаем}
\begin{itemize}
    \item Дерево проверяет совместную выполнимость уравнений (формул со знаками) $A=1$, $A=0$.
    \item Для тождественной истинности $A$ начинаем с $A=0$.
    \item Для эквивалентности строим два дерева: $A=1,B=0$ и $A=0,B=1$.
    \pause
    \item Ветвь закрыта, если на ней есть $C=1$ и $C=0$; закрытую ветвь не продолжаем.
    \item Если все ветви закрыты, исходные уравнения невыполнимы.
    \item Сначала рассматриваем замкнутые формулы без равенства, константных и функциональных символов.
\end{itemize}
\end{frame}

\begin{frame}
\frametitle{Правила деревьев истинности}
\small
\setlength{\parskip}{0pt}
\begin{center}
    $\dfrac{\neg A=1}{A=0},\qquad \dfrac{\neg A=0}{A=1}$.

    \medskip
    \begin{tabular}{|c|c|c||c|c|c|}
        \hline
        $\alpha$ & $\alpha_1$ & $\alpha_2$ & $\beta$ & $\beta_1$ & $\beta_2$ \\
        \hline
        $A\land B=1$ & $A=1$ & $B=1$ & $A\land B=0$ & $A=0$ & $B=0$ \\
        $A\lor B=0$ & $A=0$ & $B=0$ & $A\lor B=1$ & $A=1$ & $B=1$ \\
        $A\to B=0$ & $A=1$ & $B=0$ & $A\to B=1$ & $A=0$ & $B=1$ \\
        \hline
    \end{tabular}

    \smallskip
    $\alpha$: выполняются оба уравнения; $\beta$: хотя бы одно.\\
    Для $\alpha$: обе строки в одной ветви; для $\beta$: две ветви.
    \pause

    \medskip
    \begin{tabular}{|c|c||c|c|}
        \hline
        $\gamma$ & $\gamma(t)$ & $\delta$ & $\delta(a)$ \\
        \hline
        $\forall v\,A(v)=1$ & $A(t)=1$ & $\forall v\,A(v)=0$ & $A(a)=0$ \\
        $\exists v\,A(v)=0$ & $A(t)=0$ & $\exists v\,A(v)=1$ & $A(a)=1$ \\
        \hline
    \end{tabular}

    \smallskip
    $\gamma$: $t$ --- замкнутый терм; правило можно повторять.\\
    $\delta$: $a$ --- новый для ветви параметр.

    \smallskip
    Раскрывая вершину, продолжаем \textbf{все открытые ветви},\\
    проходящие через неё; результаты дописываем в их конец.
\end{center}
\end{frame}

\begin{frame}
\frametitle{Деревья без знаков: напоминание}
\begin{itemize}
    \item Вместо формул со знаками можно писать только формулы:
    \[
        A=1\ \rightsquigarrow\ A,
        \qquad A=0\ \rightsquigarrow\ \neg A.
    \]
    \item Ветвь закрывается, если содержит $A$ и $\neg A$.
    \pause
    \item Правила для кванторов:
\end{itemize}
\begin{center}
    \begin{tabular}{|c|c||c|c|}
        \hline
        $\gamma$ & $\gamma(t)$ & $\delta$ & $\delta(a)$ \\
        \hline
        $\forall x\,A(x)$ & $A(t)$ & $\exists x\,A(x)$ & $A(a)$ \\
        $\neg\exists x\,A(x)$ & $\neg A(t)$ & $\neg\forall x\,A(x)$ & $\neg A(a)$ \\
        \hline
    \end{tabular}
\end{center}
\begin{itemize}
    \item Условия те же: $t$ --- замкнутый терм, $a$ --- новый для ветви параметр; $\gamma$ можно раскрывать повторно.
    \item Правила для связок те же; например $\dfrac{\neg\neg A}{A}$.
\end{itemize}
\end{frame}

\begin{frame}
\frametitle{Параметры и повторное раскрытие}
\begin{itemize}
    \item Параметры $a_1,a_2,\ldots$ обозначают объекты модели.
    \item $A(t)$ означает подстановку терма $t$ вместо свободных вхождений переменной $v$ в $A(v)$.
    \item В наших примерах замкнутые термы --- параметры.
    \pause
    \item Для $\delta$ берём параметр, ещё не встречавшийся на ветви. Разные имена не требуют разных объектов.
    \item Для $\gamma$ записываем использованные параметры, но не ставим $\checkmark$: могут понадобиться новые подстановки.
    \pause
    \item Если параметров ещё нет, вводим $a_1$: универсум непуст.
    \item Запись $\forall x\,A(x)=0$ означает $(\forall x\,A(x))=0$, не $\forall x\,(A(x)=0)$.
    \item Без знаков различие видно явно: $\neg\forall x\,A(x)$ и $\forall x\,\neg A(x)$.
\end{itemize}
\end{frame}

\begin{frame}
\frametitle{Пример дерева истинности}
\begin{itemize}
    \item Первый пример: одно из правил де Моргана для кванторов, $\neg (\forall x ~ P(x)) \to (\exists x ~ \neg P(x))$.
    \item[] 
\only<+>{
    \begin{tableau}{
    }
    [{\neg (\forall x ~ P(x)) \to (\exists x ~ \neg P(x)) = 0}, just=Дано]
    \end{tableau}
}
\only<+>{
    \begin{tableau}{
    }
    [{\neg (\forall x ~ P(x)) \to (\exists x ~ \neg P(x)) = 0}, just=Дано, checked
    [{\neg (\forall x ~ P(x)) = 1}, just=1, checked
    [{\exists x ~ \neg P(x) = 0}, just = 1
    [{\forall x ~ P(x) = 0}, just=2
    ]]]]
    \end{tableau}
}
\only<+>{
    \begin{tableau}{
        }
        [{\neg (\forall x ~ P(x)) \to (\exists x ~ \neg P(x)) = 0}, just=Дано, checked
        [{\neg (\forall x ~ P(x)) = 1}, just=1, checked
        [{\exists x ~ \neg P(x) = 0}, just = 1
        [{\forall x ~ P(x) = 0}, just=2, checked=a_1
        [{P(a_1) = 0}, just=4, checked
        ]]]]]
    \end{tableau}
}
\only<+->{
    \begin{tableau}{
        }
        [{\neg (\forall x ~ P(x)) \to (\exists x ~ \neg P(x)) = 0}, just=Дано, checked
        [{\neg (\forall x ~ P(x)) = 1}, just=1, checked
        [{\exists x ~ \neg P(x) = 0}, just = 1, subs=a_1
        [{\forall x ~ P(x) = 0}, just=2, checked=a_1
        [{P(a_1) = 0}, just=4, checked
        [{\neg P(a_1) = 0}, just=3, checked
        [{P(a_1) = 1}, just=6, checked, close={5,7}
        ]]]]]]]
    \end{tableau}
}
\item<5-> Дерево закрылось, значит формула $\neg (\forall x ~ P(x)) \to (\exists x ~ \neg P(x))$ тождественно истинна.
\end{itemize}
\end{frame}

\begin{frame}
\frametitle{Порядок применения правил}
\begin{itemize}
    \item Раскрывая вершину, дописываем результат в конец \textbf{всех открытых ветвей}, проходящих через неё.
    \item Если сначала применить $\gamma$ к $a_1$, правило $\delta$ потребует новый $a_2$. Придётся повторить $\gamma$ с $a_2$.
    \pause
    \item Обычно удобно сначала раскрывать $\neg$, $\alpha$ и $\delta$, затем $\beta$, затем $\gamma$.
    \item После появления новых параметров возвращаемся к $\gamma$-формулам.
    \pause
    \item Не откладываем другие формулы, бесконечно раскрывая одну $\gamma$-формулу. При бесконечном построении нужна систематическая процедура.
    \item При ветвлении следим за подстановками и свежестью параметров отдельно в каждой ветви.
\end{itemize}
\end{frame}




\begin{frame}
\begin{framefont}{\small}
\frametitle{Пример дерева истинности (2)}
\begin{itemize}
    \item Проверим $F=(\forall x\,P(x)\land Q(x))\to(\forall x\,P(x))\land(\forall x\,Q(x))$.
    \item[] 
    \only<+>{
        \begin{tableau}{}
        [{F=0}, just=Дано, checked
[{(\forall x\,P(x)\land Q(x))=1}, just=1
            [{(\forall x ~ P(x)) \land (\forall x ~ Q(x)) = 0}, just=1]]
        ]
    \end{tableau}
    }
    \only<+>{
    \begin{tableau}{}
        [{F=0}, just=Дано, checked
[{(\forall x\,P(x)\land Q(x))=1}, just=1
        [{(\forall x ~ P(x)) \land (\forall x ~ Q(x)) = 0}, just=1, checked
        [{\forall x ~ P(x) = 0}, just = 3
        ]
        [{\forall x ~ Q(x) = 0}, just = 3
        ]]]
        ]
    \end{tableau}
    }
    \only<+>{
        \begin{tableau}{}
        [{F=0}, just=Дано, checked
[{(\forall x\,P(x)\land Q(x))=1}, just=1
        [{(\forall x ~ P(x)) \land (\forall x ~ Q(x)) = 0}, just=1, checked
        [{\forall x ~ P(x) = 0}, just = 3, checked=a_1
        [{P(a_1) = 0}, just=4, checked
        ]]
        [{\forall x ~ Q(x) = 0}, just = 3, checked=a_1
        [{Q(a_1) = 0}, just=4, checked
        ]
        ]]]
        ]
    \end{tableau}
    }
    \only<+>{
        \begin{tableau}{}
        [{F=0}, just=Дано, checked
[{(\forall x\,P(x)\land Q(x))=1}, just=1, subs=a_1
        [{(\forall x ~ P(x)) \land (\forall x ~ Q(x)) = 0}, just=1, checked
        [{\forall x ~ P(x) = 0}, just = 3, checked=a_1
        [{P(a_1) = 0}, just=4, checked
        [{P(a_1) \land Q(a_1) = 1}, just=2
        ]]]
        [{\forall x ~ Q(x) = 0}, just = 3, checked=a_1
        [{Q(a_1) = 0}, just=4, checked
        [{P(a_1) \land Q(a_1) = 1}, just=2
        ]]
        ]]]
        ]
    \end{tableau}
    }
    \only<+->{
    \begin{tableau}{}
        [{F=0}, just=Дано, checked
[{(\forall x\,P(x)\land Q(x))=1}, just=1, subs=a_1
        [{(\forall x ~ P(x)) \land (\forall x ~ Q(x)) = 0}, just=1, checked
        [{\forall x ~ P(x) = 0}, just=3, checked=a_1
        [{P(a_1) = 0}, just=4, checked
        [{P(a_1) \land Q(a_1) = 1}, just=2, checked
        [{P(a_1) = 1}, just=6, checked
        [{Q(a_1) = 1}, just=6, checked, close={5,7}
        ]]]]]
        [{\forall x ~ Q(x) = 0}, just=3, checked=a_1
        [{Q(a_1) = 0}, just=4, checked
        [{P(a_1) \land Q(a_1) = 1}, just=2, checked
        [{P(a_1) = 1}, just=6, checked
        [{Q(a_1) = 1}, just=6, checked, close={5,8}
        ]]]]
        ]]]
        ]
    \end{tableau}
}
    \item<6-> Дерево закрылось, значит формула тождественно истинна.
    \item<3-> В строке 4 можем использовать $a_1$ в обеих ветвях.
    \item<4-> Строку 2 раскрываем в обеих открытых ветвях под ней.
\end{itemize}
\end{framefont}
\end{frame}

\begin{frame}
\begin{framefont}{\footnotesize}
\frametitle{Пример дерева истинности (3)}
\begin{itemize}
    \item $F=(\forall x\,P(x)\lor Q(x))\to(\forall x\,P(x))\lor(\forall x\,Q(x))$. \\ Тождественно истинна ли эта формула?
    
    \only<+>{
        \begin{tableau}{}
        [{F=0}, just=Дано, checked
[{(\forall x\,P(x)\lor Q(x))=1}, just=1
            [{(\forall x ~ P(x)) \lor (\forall x ~ Q(x)) = 0}, just=1
            ]]
        ]
    \end{tableau}
    }
    
    \only<+>{
        \begin{tableau}{}
        [{F=0}, just=Дано, checked
[{(\forall x\,P(x)\lor Q(x))=1}, just=1
            [{(\forall x ~ P(x)) \lor (\forall x ~ Q(x)) = 0}, just=1, checked
            [{\forall x ~ P(x) = 0}, just = 3
            [{\forall x ~ Q(x) = 0}, just = 3
            ]]]]
        ]
    \end{tableau}
    }
    
    \only<+>{
        \begin{tableau}{}
        [{F=0}, just=Дано, checked
[{(\forall x\,P(x)\lor Q(x))=1}, just=1
            [{(\forall x ~ P(x)) \lor (\forall x ~ Q(x)) = 0}, just=1, checked
            [{\forall x ~ P(x) = 0}, just = 3, checked=a_1
            [{\forall x ~ Q(x) = 0}, just = 3, checked=a_2
            [{P(a_1) = 0}, just=4, checked
            [{Q(a_2) = 0}, just=5, checked
            ]]]]]]
        ]
    \end{tableau}
    }
    
        \only<+>{
        \begin{tableau}{}
        [{F=0}, just=Дано, checked
[{(\forall x\,P(x)\lor Q(x))=1}, just=1, subs={a_1}
            [{(\forall x ~ P(x)) \lor (\forall x ~ Q(x)) = 0}, just=1, checked
            [{\forall x ~ P(x) = 0}, just = 3, checked=a_1
            [{\forall x ~ Q(x) = 0}, just = 3, checked=a_2
            [{P(a_1) = 0}, just=4, checked
            [{Q(a_2) = 0}, just=5, checked
            [{P(a_1) \lor Q(a_1) = 1}, just=2
            ]]]]]]]
        ]
    \end{tableau}
    }

    \only<+>{
    \begin{tableau}{}
        [{F=0}, just=Дано, checked
[{(\forall x\,P(x)\lor Q(x))=1}, just=1, subs={a_1}
        [{(\forall x ~ P(x)) \lor (\forall x ~ Q(x)) = 0}, just=1, checked
        [{\forall x ~ P(x) = 0}, just = 3, checked=a_1
        [{\forall x ~ Q(x) = 0}, just = 3, checked=a_2
        [{P(a_1) = 0}, just=4, checked
        [{Q(a_2) = 0}, just=5, checked
        [{P(a_1) \lor Q(a_1) = 1}, just=2, checked
        [{P(a_1) = 1}, just={8 (ветвь закончена?)}, checked, close={6,9}
        ]
        [{Q(a_1) = 1}, just={8 (ветвь закончена?)}, checked
        ]]]]]]]]
        ]
    \end{tableau}
}

    \only<+>{
        \begin{tableau}{}
        [{F=0}, just=Дано, checked
[{(\forall x\,P(x)\lor Q(x))=1}, just=1, subs={a_1,a_2}
[{(\forall x ~ P(x)) \lor (\forall x ~ Q(x)) = 0}, just=1, checked
[{\forall x ~ P(x) = 0}, just = 3, checked=a_1
[{\forall x ~ Q(x) = 0}, just = 3, checked=a_2
[{P(a_1) = 0}, just=4, checked
[{Q(a_2) = 0}, just=5, checked
[{P(a_1) \lor Q(a_1) = 1}, just=2, checked
[{P(a_1) = 1}, just=8, checked, close={6,9}
]
[{Q(a_1) = 1}, just=8, checked
[{P(a_2) \lor Q(a_2) = 1}, just=2
]]]]]]]]]
        ]
    \end{tableau}
    }

    \only<+->{
    \begin{tableau}{}
        [{F=0}, just=Дано, checked
[{(\forall x\,P(x)\lor Q(x))=1}, just=1, subs={a_1,a_2}
        [{(\forall x ~ P(x)) \lor (\forall x ~ Q(x)) = 0}, just=1, checked
        [{\forall x ~ P(x) = 0}, just = 3, checked=a_1
        [{\forall x ~ Q(x) = 0}, just = 3, checked=a_2
        [{P(a_1) = 0}, just=4, checked
        [{Q(a_2) = 0}, just=5, checked
        [{P(a_1) \lor Q(a_1) = 1}, just=2, checked
        [{P(a_1) = 1}, just=8, checked, close={6,9}
        ]
        [{Q(a_1) = 1}, just=8, checked
        [{P(a_2) \lor Q(a_2) = 1}, just=2, checked
        [{P(a_2) = 1}, just=10, checked
        ]
        [{Q(a_2) = 1}, just=10, checked, close={7,11}
        ]
        ]]]]]]]]]
        ]
    \end{tableau}
}
    \item<7-> Дерево не закрылось!
\end{itemize}
\end{framefont}
\end{frame}

\begin{frame}
\frametitle{Контрмодель из законченной ветви}
\begin{itemize}
    \item Ветвь закончена: раскрыты все $\neg$-, $\alpha$-, $\beta$- и $\delta$-формулы, а все $\gamma$ --- со всеми параметрами ветви (хотя бы с одним).
    \pause
    \item Возьмём универсум $\{a_1,a_2\}$; атомы ветви задают таблицу:
    \item[] \begin{tabular}{|c|cc|}
        \hline
        $x$ & $a_1$ & $a_2$ \\ \hline
        $P(x)$ & $0$ & $1$ \\ \hline
        $Q(x)$ & $1$ & $0$ \\ \hline
    \end{tabular}
    \pause
    \item Если какого-то атома в законченной ветви нет, его значение можно выбрать произвольно.
    \pause
    \item Проверим: $\forall x\,(P(x)\lor Q(x))$ истинна.
    \item Но $\forall x\,P(x)$ и $\forall x\,Q(x)$ ложны. Значит, заключение ложно и это контрмодель.
\end{itemize}
\end{frame}

\begin{frame}
\begin{framefont}{\small}
\frametitle{Пример дерева истинности (4)}
\begin{itemize}
    \item $F=(\forall x\,\exists y\,P(x,y))\to(\exists y\,\forall x\,P(x,y))$. \\ Тождественно истинна ли эта формула?
    \item[] 
\only<+>{
    \begin{tableau}{}
        [{F=0}, just=Дано, checked
[{\forall x \exists y ~ P(x, y) = 1}, just=1
        [{\exists y \forall x ~ P(x, y) = 0}, just=1
        ]]
        ]
    \end{tableau}
}

    \only<+>{
    \begin{tableau}{}
        [{F=0}, just=Дано, checked
[{\forall x \exists y ~ P(x, y) = 1}, just=1, subs={a_1}
        [{\exists y \forall x ~ P(x, y) = 0}, just=1
        [{\exists y ~ P(a_1, y) = 1}, just=2, checked=a_2
        [{P(a_1, a_2) = 1}, just=4, checked
        ]]]]
        ]
    \end{tableau}
}

    \only<+>{
\begin{tableau}{}
        [{F=0}, just=Дано, checked
[{\forall x \exists y ~ P(x, y) = 1}, just=1, subs={a_1}
    [{\exists y \forall x ~ P(x, y) = 0}, just=1, subs={a_2}
    [{\exists y ~ P(a_1, y) = 1}, just=2, checked=a_2
    [{P(a_1, a_2) = 1}, just=4, checked
    [{\forall x ~ P(x, a_2) = 0}, just=3, checked=a_3
    [{P(a_3, a_2) = 0}, just={6 (ветвь закончена?)}, checked
    ]]]]]]
        ]
    \end{tableau}
}

    \only<+->{
    \begin{tableau}{}
        [{F=0}, just=Дано, checked
[{\forall x \exists y ~ P(x, y) = 1}, just=1, subs={a_1,a_3,\ldots}
        [{\exists y \forall x ~ P(x, y) = 0}, just=1, subs={a_2,\ldots}
        [{\exists y ~ P(a_1, y) = 1}, just=2, checked=a_2
        [{P(a_1, a_2) = 1}, just=4, checked
        [{\forall x ~ P(x, a_2) = 0}, just=3, checked=a_3
        [{P(a_3, a_2) = 0}, just=6, checked
        [{\exists y ~ P(a_3, y) = 1}, just=2, checked=a_4
        [{\ldots}
        ]]]]]]]]
        ]
    \end{tableau}
}
    \onslide<+->
    \item После любого конечного числа шагов ветвь открыта, но ещё не закончена: появляются новые параметры и подстановки.
    \onslide<+->
    \item Чтобы опровергнуть формулу, построим контрмодель.
\end{itemize}
\end{framefont}
\end{frame}

\begin{frame}
\begin{framefont}{\small}
\frametitle{Условия на контрмодель}
\begin{itemize}
    \item В нашей ветви видны некоторые значения предиката:
    \item[] \begin{tabular}{c | c c c c c}
        \diagbox[height=1.5\line]{x}{y} & $a_1$ & $a_2$ & $a_3$ & $a_4$ & \ldots \\ \hline
        $a_1$ & ? & 1 & ? & ? & \\
        $a_2$ & ? & ? & ? & ? & \\
        $a_3$ & ? & 0 & ? & 1 & \\
        $a_4$ & ? & ? & ? & ? & \\
        \ldots & & & & & \\ 
    \end{tabular}
    \item Как в терминах этой таблицы выглядят:
    \begin{itemize}
        \item $\forall x \exists y ~ P(x, y) = 1$? \pause В каждой строке есть 1. \pause
        \item $\exists y \forall x ~ P(x, y) = 0$? \pause Нет столбца, где все 1. 
        \item[] Помните, что это $(\exists y \forall x ~ P(x, y)) = 0$, а не $\exists y \forall x ~ (P(x, y) = 0)$!
    \end{itemize}
\end{itemize}
\end{framefont}
\end{frame}

\begin{frame}
\begin{framefont}{\small}
\frametitle{Бесконечная и конечная контрмодели}
\begin{itemize}
    \item Расставить значения произвольно нельзя. Возьмём $P(a_i,a_j)=1$ ровно при $j=i+1$:
    \item[] {\small \begin{tabular}{c | c c c c c}
        \diagbox[height=1.5\line]{x}{y} & $a_1$ & $a_2$ & $a_3$ & $a_4$ & \ldots \\ \hline
        $a_1$ & 0 & 1 & 0 & 0 & \\
        $a_2$ & 0 & 0 & 1 & 0 & \\
        $a_3$ & 0 & 0 & 0 & 1 & \\
        $a_4$ & 0 & 0 & 0 & 0 & \\
        \ldots & & & & & \\ 
    \end{tabular}}
    \pause
    \item А есть ли конечная контрмодель? Условия те же:
    \item В каждой строке есть 1; нет столбца, где все 1. 
    \pause
    \item Конечно, есть. Например:
    \item[] {\small \begin{tabular}{c|cc}
        \diagbox[height=1.5\line]{x}{y} & $a_1$ & $a_2$ \\ \hline
        $a_1$ & 1 & 0 \\
        $a_2$ & 0 & 1 \\
        \hline
    \end{tabular}}
\end{itemize}
\end{framefont}
\end{frame}

\begin{frame}
\frametitle{Когда можно остановиться}
\begin{itemize}
    \item Закрытое дерево доказывает невыполнимость исходных уравнений.
    \item Для контрмодели нужна законченная открытая ветвь:
    \begin{itemize}
        \item раскрыты все $\neg$-, $\alpha$-, $\beta$- и $\delta$-формулы;
        \item каждая $\gamma$-формула раскрыта со всеми параметрами ветви, хотя бы с одним.
    \end{itemize}
    \pause
    \item В наших примерах параметры такой ветви можно взять за объекты модели; атомы задают значения предикатов.
    \item Проверяем в построенной модели все посылки и ложность заключения.
    \pause
    \item Просто открытая ветвь ещё не даёт контрмодель. Построение может продолжаться бесконечно, даже если есть конечный контрпример.
\end{itemize}
\end{frame}

\begin{frame}
\frametitle{Константы и функциональные символы}
\begin{itemize}
    \item Теперь разрешим в сигнатуре константы и функции. Замкнутые термы: $c$, $a_1$, $f(a_1)$, $g(c,f(a_2))$,\ldots
    \item В правиле $\gamma$ можно использовать любой такой терм $t$:
    \[
        \dfrac{\forall x\,A(x)=1}{A(t)=1},
        \qquad\text{например:}\quad
        \dfrac{\forall x\,P(x)=1}{P(f(a_1))=1}.
    \]
    \pause
    \item В правиле $\delta$ по-прежнему нужен новый параметр, а не произвольный терм.
    \item Отдельного правила раскрытия $f(t)$ нет: это терм, а не формула.
    \pause
    \item Для законченной ветви $\gamma$-формулы должны быть раскрыты со всеми замкнутыми термами из символов сигнатуры и параметров ветви, даже если терм ещё не записан в ветви.
\end{itemize}
\end{frame}

\begin{frame}
\frametitle{Пример с функциональным символом}
\begin{itemize}
    \item Проверим $F=(\forall x\,P(x))\to(\forall x\,P(f(x)))$.
    \item[]
    \only<1>{
        \begin{tableau}{}
        [{F=0}, just=Дано, checked
[{\forall x\,P(x)=1}, just=1
            [{\forall x\,P(f(x))=0}, just=1]]
        ]
    \end{tableau}
    }
    \only<2>{
        \begin{tableau}{}
        [{F=0}, just=Дано, checked
[{\forall x\,P(x)=1}, just=1
            [{\forall x\,P(f(x))=0}, just=1, checked=a_1
            [{P(f(a_1))=0}, just=3, checked]]]
        ]
    \end{tableau}
    }
    \only<3->{
        \begin{tableau}{}
        [{F=0}, just=Дано, checked
[{\forall x\,P(x)=1}, just=1, subs={f(a_1)}
            [{\forall x\,P(f(x))=0}, just=1, checked=a_1
            [{P(f(a_1))=0}, just=3, checked
            [{P(f(a_1))=1}, just=2, close={4,5}, checked]]]]
        ]
    \end{tableau}
    }
    \item<2-> $\delta$ вводит $a_1$, а $\gamma$ применяем к терму $f(a_1)$.
    \item<3-> Дерево закрыто: формула тождественно истинна.
\end{itemize}
\end{frame}

\begin{frame}
\begin{framefont}{\small}
\frametitle{Повторное раскрытие с разными термами}
\begin{itemize}
    \item $F=(\forall x\,P(x)\to P(f(x)))\to(\forall x\,P(x)\to P(f(f(x))))$.
    \item[]
\only<1>{
{\footnotesize
\begin{tableau}{}
[{F=0}, just=Дано, checked
  [{\forall x\,(P(x)\to P(f(x)))=1}, just=1
    [{\forall x\,(P(x)\to P(f(f(x))))=0}, just=1]
  ]
]
\end{tableau}
}
}
\only<2>{
{\footnotesize
\begin{tableau}{}
[{F=0}, just=Дано, checked
  [{\forall x\,(P(x)\to P(f(x)))=1}, just=1
    [{\forall x\,(P(x)\to P(f(f(x))))=0}, just=1, checked=a_1
      [{P(a_1)\to P(f(f(a_1)))=0}, just=3, checked
        [{P(a_1)=1}, just=4, checked
          [{P(f(f(a_1)))=0}, just=4, checked]
        ]
      ]
    ]
  ]
]
\end{tableau}
}
}
\only<3>{
{\footnotesize
\begin{tableau}{}
[{F=0}, just=Дано, checked
  [{\forall x\,(P(x)\to P(f(x)))=1}, just=1, subs={a_1}
    [{\forall x\,(P(x)\to P(f(f(x))))=0}, just=1, checked=a_1
      [{P(a_1)\to P(f(f(a_1)))=0}, just=3, checked
        [{P(a_1)=1}, just=4, checked
          [{P(f(f(a_1)))=0}, just=4, checked
            [{P(a_1)\to P(f(a_1))=1}, just=2, checked
              [{P(a_1)=0}, just=7, close={5,8}, checked]
              [{P(f(a_1))=1}, just=7, checked]
            ]
          ]
        ]
      ]
    ]
  ]
]
\end{tableau}
}
}
\only<4>{
{\footnotesize
\begin{tableau}{}
[{F=0}, just=Дано, checked
  [{\forall x\,(P(x)\to P(f(x)))=1}, just=1, subs={a_1,f(a_1)}
    [{\forall x\,(P(x)\to P(f(f(x))))=0}, just=1, checked=a_1
      [{P(a_1)\to P(f(f(a_1)))=0}, just=3, checked
        [{P(a_1)=1}, just=4, checked
          [{P(f(f(a_1)))=0}, just=4, checked
            [{P(a_1)\to P(f(a_1))=1}, just=2, checked
              [{P(a_1)=0}, just=7, close={5,8}, checked]
              [{P(f(a_1))=1}, just=7, checked
                [{P(f(a_1))\to P(f(f(a_1)))=1}, just=2]
              ]
            ]
          ]
        ]
      ]
    ]
  ]
]
\end{tableau}
}
}
\only<5->{
{\footnotesize
\begin{tableau}{}
[{F=0}, just=Дано, checked
  [{\forall x\,(P(x)\to P(f(x)))=1}, just=1, subs={a_1,f(a_1)}
    [{\forall x\,(P(x)\to P(f(f(x))))=0}, just=1, checked=a_1
      [{P(a_1)\to P(f(f(a_1)))=0}, just=3, checked
        [{P(a_1)=1}, just=4, checked
          [{P(f(f(a_1)))=0}, just=4, checked
            [{P(a_1)\to P(f(a_1))=1}, just=2, checked
              [{P(a_1)=0}, just=7, close={5,8}, checked]
              [{P(f(a_1))=1}, just=7, checked
                [{P(f(a_1))\to P(f(f(a_1)))=1}, just=2, checked
                  [{P(f(a_1))=0}, just=9, close={8,10}, checked]
                  [{P(f(f(a_1)))=1}, just=9, close={6,10}, checked]
                ]
              ]
            ]
          ]
        ]
      ]
    ]
  ]
]
\end{tableau}
}
}
    \item<4-> Строку 2 раскрыли сначала с $a_1$, затем с $f(a_1)$.
    \item<5-> Все ветви закрыты: формула тождественно истинна.
\end{itemize}
\end{framefont}
\end{frame}

% Правила равенства: https://builds.openlogicproject.org/content/first-order-logic/tableaux/identity.pdf
\begin{frame}
\frametitle{Правила для равенства}
\begin{itemize}
    \item $s$ и $t$ --- замкнутые термы. Равенство обозначает совпадение объектов, а не произвольный предикат.
    \item Рефлексивность и замена равных:
    \[
        \dfrac{}{(t=t)=1}
        \qquad
        \dfrac{(s=t)=1\quad A[s]=i}{A[t]=i},\quad i\in\{0,1\}.
    \]
    \item $A[t]$ получается из $A[s]$ заменой выбранных вхождений терма $s$ на $t$; можно заменять и в обратную сторону.
    \pause
    \item Заменять можно и внутри термов: например, $P(f(s))=0$ даёт $P(f(t))=0$ при $(s=t)=1$.
    \item Атомы можно повторно использовать в правилах равенства даже с $\checkmark$: раскрывать в них нечего.
    \pause
    \item Ветвь с $(t=t)=0$ сразу закрывается. Из $(s=t)=0$ правило замены применять нельзя.
\end{itemize}
\end{frame}

\begin{frame}
\begin{framefont}{\small}
\frametitle{Правила Ривза (Reeves): предикаты}
% https://staff.science.uva.nl/u.endriss/teaching/ilcs/2006/slides/more-fol-tableaux-4up.pdf
\begin{itemize}
    \item $t_1,\ldots,t_n,s_1,\ldots,s_n$ --- замкнутые термы, $n\ge1$. Две исходные формулы должны быть на одной ветви.
    \begin{center}
        \begin{tableau}{not line numbering, for tree={l sep=3mm}}
            [{P(t_1,\ldots,t_n)=1}
            [{P(s_1,\ldots,s_n)=0}
                [{(t_1=s_1)=0}]
                [{\cdots}]
                [{(t_n=s_n)=0}]
            ]]
        \end{tableau}
    \end{center}
    \pause
    \item Если значения одного предиката различны, хотя бы одна пара соответствующих аргументов различна.
    \item Сразу создаём $n$ ветвей: в ветви $i$ добавляем $(t_i=s_i)=0$.
    \item $P$ здесь может быть и предикатом равенства.
    \pause
    \item Замена равных может породить бесконечно много формул: при $f(a)=a$ из $P(a)$ получаем $P(f(a))$, $P(f(f(a)))$,\ldots
    \item Ривз работает с атомами одного предиката с противоположными знаками: меньше лишних замен, но ветвление может увеличить дерево.
\end{itemize}
\end{framefont}
\end{frame}

\begin{frame}
\begin{framefont}{\small}
\frametitle{Правила Ривза: функции и равенство}
\begin{itemize}
    \item Для одного и того же $n$-местного функционального символа:
    \begin{center}
        \begin{tableau}{not line numbering, for tree={l sep=3mm}}
            [{(f(t_1,\ldots,t_n)=f(s_1,\ldots,s_n))=0}
                [{(t_1=s_1)=0}]
                [{\cdots}]
                [{(t_n=s_n)=0}]
            ]
        \end{tableau}
    \end{center}
    \item Разные результаты требуют хотя бы одной пары разных аргументов. Для $n=1$ получаем $(t_1=s_1)=0$ без ветвления.
    \pause
    \item Дополняем правила симметричностью:
    \[ \dfrac{(s=t)=1}{(t=s)=1}. \]
    \item Сохраняем обычное закрытие ветви и закрытие при $(t=t)=0$.
    \pause
    \item Эти правила можно использовать вместо замены равных. Применения можно повторять для новых пар атомов и новых неравенств.
\end{itemize}
\end{framefont}
\end{frame}

\begin{frame}
\frametitle{Пример: равенство аргументов функции}
\begin{itemize}
    \item Проверим $F=(c=d)\to(f(c)=f(d))$.
    \item[]
    \only<1>{
        \begin{tableau}{}
        [{F=0}, just=Дано, checked
[{(c=d)=1}, just=1, checked
            [{(f(c)=f(d))=0}, just=1, checked]]
        ]
    \end{tableau}
    }
    \only<2>{
        \begin{tableau}{}
        [{F=0}, just=Дано, checked
[{(c=d)=1}, just=1, checked
            [{(f(c)=f(d))=0}, just=1, checked
            [{(f(c)=f(c))=1}, just=рефл., checked]]]
        ]
    \end{tableau}
    }
    \only<3->{
        \begin{tableau}{}
        [{F=0}, just=Дано, checked
[{(c=d)=1}, just=1, checked
            [{(f(c)=f(d))=0}, just=1, checked
            [{(f(c)=f(c))=1}, just=рефл., checked
            [{(f(c)=f(d))=1}, just={=, 2, 4}, close={3,5}, checked]]]]
        ]
    \end{tableau}
    }
    \item<3-> В строке 4 заменили $c$ на $d$ только в правой части. Дерево закрыто.
    \item<4-> Обратная импликация неверна без инъективности $f$: разные аргументы могут иметь одинаковые значения.
\end{itemize}
\end{frame}

\begin{frame}
\begin{framefont}{\small}
\frametitle{Тот же пример по правилам Ривза}
\begin{itemize}
    \item Проверим $F=(c=d)\to(f(c)=f(d))$.
    \item[]
    \only<1>{
        \begin{tableau}{}
            [{F=0}, just=Дано, checked
            [{(c=d)=1}, just=1, checked
            [{(f(c)=f(d))=0}, just=1, checked]]]
        \end{tableau}
    }
    \only<2->{
        \begin{tableau}{}
            [{F=0}, just=Дано, checked
            [{(c=d)=1}, just=1, checked
            [{(f(c)=f(d))=0}, just=1, checked
            [{(c=d)=0}, just={Ривз, 3}, close={2,4}, checked]]]]
        \end{tableau}
    }
    \item<2-> Разные значения $f(c)$ и $f(d)$ требуют $c\ne d$, что противоречит строке 2. Ветвь закрыта.
    \item<2-> Для этого примера достаточно одного применения нового правила; добавлять $f(c)=f(c)$ не понадобилось.
\end{itemize}
\end{framefont}
\end{frame}

\begin{frame}
\begin{framefont}{\small}
\frametitle{Контрмодели с функциями и равенством}
\begin{itemize}
    \item Для законченной ветви учитываем все замкнутые термы и все нужные следствия правил равенства.
    \item При построении модели берём термы за имена объектов. Термы, равенство которых следует из ветви, объединяем в один класс $[t]$.
    \pause
    \item Константа $c$ обозначает $[c]$, а функцию задаём так:
    \[ f^M([t_1],\ldots,[t_n])=[f(t_1,\ldots,t_n)]. \]
    \item Равные аргументы должны давать одинаковые значения функций и предикатов; $(s=t)=0$ требует разных классов.
    \pause
    \item Замкнутых термов может быть бесконечно много: $a$, $f(a)$, $f(f(a))$,\ldots
    \item Иногда проще предложить конечную модель напрямую. Проверяем в ней все исходные уравнения, включая равенства и неравенства.
\end{itemize}
\end{framefont}
\end{frame}

\begin{frame}
\begin{framefont}{\small}
\frametitle{Обратная импликация: дерево}
\begin{itemize}
    \item Проверим обратную к формуле из примера с функциональным символом:
    \[ F=(\forall x\,P(f(x)))\to(\forall x\,P(x)). \]
    \item[]
    \only<1>{
        \begin{tableau}{}
            [{F=0}, just=Дано, checked
            [{\forall x\,P(f(x))=1}, just=1
            [{\forall x\,P(x)=0}, just=1]]]
        \end{tableau}
    }
    \only<2>{
        \begin{tableau}{}
            [{F=0}, just=Дано, checked
            [{\forall x\,P(f(x))=1}, just=1
            [{\forall x\,P(x)=0}, just=1, checked=a_1
            [{P(a_1)=0}, just=3, checked]]]]
        \end{tableau}
    }
    \only<3->{
        \begin{tableau}{}
            [{F=0}, just=Дано, checked
            [{\forall x\,P(f(x))=1}, just=1, subs={a_1,f(a_1),\ldots}
            [{\forall x\,P(x)=0}, just=1, checked=a_1
            [{P(a_1)=0}, just=3, checked
            [{P(f(a_1))=1}, just=2, checked
            [{P(f(f(a_1)))=1}, just=2, checked
            [{\ldots}]]]]]]]
        \end{tableau}
    }
    \item<3-> Ветвь открыта, но ещё не закончена: есть всё новые термы для раскрытия строки 2.
    \item<3-> Предложим конечную модель и проверим в ней $F=0$.
\end{itemize}
\end{framefont}
\end{frame}

\begin{frame}
\frametitle{Обратная импликация: контрмодель}
\begin{itemize}
    \item Возьмём универсум $\{a,b\}$:
    \item[] {\setlength{\tabcolsep}{1em}
    \begin{tabular}{c|c|c}
        $x$ & $a$ & $b$ \\ \hline
        $f(x)$ & $b$ & $b$ \\
        $P(x)$ & $0$ & $1$
    \end{tabular}}
    \pause
    \item Для каждого $x$ имеем $f(x)=b$ и $P(b)=1$. Поэтому $\forall x\,P(f(x))$ истинна.
    \item Но $P(a)=0$, поэтому $\forall x\,P(x)$ ложна.
    \pause
    \item Значит, импликация $(\forall x\,P(f(x)))\to(\forall x\,P(x))$ не тождественно истинна.
    \item При $a_1=a$ все термы $f(a_1)$, $f(f(a_1))$,\ldots\ обозначают один объект $b$.
    \item Формула $\forall x\,P(f(x))$ проверяет $P$ только на значениях функции $f$, а они могут не охватывать весь универсум.
\end{itemize}
\end{frame}

\begin{frame}
\frametitle{Итоги}
\begin{itemize}
    \item Пропозициональные правила сохраняются; для кванторов добавляются $\gamma$ и $\delta$.
    \item $\delta$ вводит нового свидетеля; $\gamma$ проверяет формулу для выбранного объекта и может применяться снова.
    \item Новые параметры и замкнутые термы требуют новых подстановок в $\gamma$-формулы.
    \item Равенство позволяет заменять равные термы, сохраняя знак формулы; $(t=t)=0$ закрывает ветвь.
    \item Все ветви закрыты --- контрпримера нет.
    \item Законченная открытая ветвь даёт контрмодель; незаконченная требует продолжения.
    \item В отличие от логики высказываний, построение не обязано закончиться за конечное число шагов.
\end{itemize}
\end{frame}

\end{document}
